% !TeX root = tesis.tex

Entre los distintos tipos de argumentos de conocimiento introducidos en la sección anterior, los \textit{\SNARK (Succint Non-Interactive ARgument of Knowledge)} ocupan un lugar central por su eficiencia y aplicabilidad práctica. Fueron introducidos por primera vez en \cite{Bitansky2012}. Estos sistemas permiten demostrar la corrección de cómputos arbitrariamente complejos mediante pruebas extremadamente cortas y verificables en tiempo reducido. La combinación de no interactividad, sucintez y solidez convierte a los \SNARK en una herramienta fundamental en criptografía moderna, particularmente en contextos donde la verificación eficiente resulta crítica, como sistemas distribuidos, protocolos de privacidad y blockchain.

\begin{definition}[\SNARK, \zkSNARK]
	Sea $\Relation \subseteq \InstanceSpace \times \WitnessSpace$ una relación en la clase de complejidad $\NP$. Un \SNARK para $\Relation$ es una tripla de algoritmos en tiempo polinomial
	$$\mathcal{S} = (\Setup, \Prover, \Verifier),$$
	definidos como:
	\begin{itemize}
		\item $\Setup(1^\securityparam)$ es un algoritmo aleatorio que genera un conjunto de parámetros públicos $\CRS$, conocido como \textit{Common Reference String}, que es usado tanto por el prover como por el verifier. $\securityparam$ es un parámetro de seguridad.
		\item $\Prover_{\CRS}: \InstanceSpace \times \WitnessSpace \to \ProofSpace$ es un algoritmo aleatorio que produce una prueba $\pi$
		\item $\Verifier_\CRS: \InstanceSpace \times \ProofSpace \to \{0,1\}$ es un algoritmo determinístico.
	\end{itemize}
	
	La tripla de algoritmos $\mathcal{S}$ debe, además, cumplir con las siguientes propiedades:
	\begin{itemize}
		\item Completitud: Para todo $(x,w) \in R$, $\Prob[\Verifier_\CRS(x,\Prover_\CRS(x,w)) = 1] = 1$.
		\item Solidez computacional: Para todo prover probabilístico en tiempo polinomial $\Prover^\star$ que produce un argumento $\pi$ tal que
		$$\Prob\left[\Verifier_\CRS(x, \pi) = 1\right]$$
		es no negligible, existe un extractor probabilístico en tiempo polinomial con acceso a $\Prover^\star$ que devuelve (extrae) un witness $w$ tal que \((x,w) \in R\). Es por esto que hablamos de que $\pi$ es un argumento y no una prueba. En adelante, cuando hablamos de \SNARK nos referimos a argumentos de conocimiento.
		\item Sucintez: Esta propiedad se divide, a su vez, en otras subpropiedades:
		\begin{itemize}
			\item Existe un polinomio $p \in \poly{N}_0$ tal que para todo $(x,w) \in \Relation$, $|\pi| \leq p(\securityparam, |x|)$, donde $\pi = \Prover_\CRS(x,w)$. En particular, el tamaño de la prueba es independiente de $|w|$.
			\item El tiempo de ejecución de $\Verifier_\CRS$ está acotado por $q(\securityparam, |x|)$, donde $q \in \poly{N_0}$. También es independiente de $|w|$.
			\item El costo de verificación de $\pi = \Prover_\CRS(x, w)$ es sublineal con respecto al costo de verificar $(x,w) \in R$.
		\end{itemize}
		\item No interactividad: El prover $\Prover_\CRS$ genera un único mensaje $\pi$, que luego es verificado por $\Verifier_\CRS$ sin interacción adicional.
	\end{itemize}
	
		Si además $\mathcal{S}$ satisface la siguiente propiedad, decimos que es un \zkSNARK.
	\begin{itemize}
		\item Conocimiento cero: Para todo $x \in \InstanceSpace$, existe un simulador probabilístico en tiempo polinomial $\Simulator_{x,\CRS}$ tal que la distribución de las pruebas simuladas $\Simulator_{x,\CRS}(x)$ es indistinguible computacionalmente de la distribución de pruebas reales $\Prover_\CRS(x,w)$, para cualquier $w \in \WitnessSpace$ tal que $(x,w) \in \Relation$. Intuitivamente, esto implica que el verifier podría haber generado por sí mismo una prueba indistinguible de la real, sin conocer el witness.
	\end{itemize}
\end{definition}

La definición de sucintez que dimos es informal. Para acercarnos a la formalidad, consideramos toda familia de instancias $(x_n,w_n) \in \Relation$ tales que el tamaño mínimo de $w_n$ para verificar $\Relation$ con $x_n$ crece de forma no acotada, el costo de verificación es asintóticamente menor que el costo de verificar explícitamente la pertenencia de $(x_n, w_n)$ a $\Relation$. Es decir, el verifier no debe volver a ejecutar el cómputo implícito en el witness.

Recordar que, como $\Relation$ está en $\NP$, entonces es seguro que existe una máquina de Turing determinística $M_\Relation$ que decide $(x,w) \in \Relation$ en tiempo polinomial respecto a $|x| + |w|$. Lo que no necesariamente es polinomial (y justamente nos sirve que no lo sea) es el lenguaje $\Language_\Relation = \{x: \exists w \mid (x,w) \in \Relation\}$, o sea, decidir la pertenencia a la relación sin conocer el \textit{witness} $w$. Tener presente que una relación en $\NP$ representa un cómputo siempre resulta útil a la hora de aprender o reforzar estos conceptos.

Ahora, vamos a hablar más informalmente de \SNARK y \zkSNARK para ganar mayor intuición respecto a estos conceptos. En primer lugar, un \SNARK puede entenderse como un mecanismo que permite a una parte demostrar que conoce un \textit{witness} $w$ consistente con un \textit{input} público $x$, sin interacción adicional y con una prueba de tamaño pequeño. La no interactividad surge del hecho de que primero se ejecuta $\Prover_\CRS$ y luego $\Verifier_\CRS$, de modo que toda la comunicación entre el \textit{prover} y el \textit{verifier} es un único mensaje $\pi$.

La propiedad de solidez computacional asegura que ningún \textit{prover} eficiente puede convencer al \textit{verifier} de una afirmación falsa, salvo con probabilidad negligible. La noción de solidez por extracción formaliza rigurosamente la intuición presentada en la sección anterior: una prueba válida no sólo convence al \textit{verifier} de la veracidad de una afirmación, sino que además certifica que el \textit{prover} posee efectivamente el conocimiento subyacente. La existencia de un extractor eficiente implica que cualquier \textit{prover} capaz de generar pruebas aceptadas con probabilidad no negligible debe, implícitamente, contener suficiente información como para reconstruir un \textit{witness} válido. En este sentido, la prueba actúa como una evidencia criptográfica de conocimiento.

La sucintez es la característica distintiva de los \SNARK. Intuitivamente, esta propiedad establece que el costo de verificación depende únicamente del tamaño del \textit{input} público $x$, y no de la complejidad del cómputo de $\Relation$. En particular, el \textit{verifier} no vuelve a ejecutar el algoritmo implícito en el \textit{witness} $w$, sino que se limita a verificar una prueba corta de la corrección del cómputo. Este desacople entre el costo de probar y el costo de verificar es lo que permite utilizar los \SNARK para la verificación de cómputos arbitrariamente grandes.

Finalmente, los \zkSNARK tienen la propiedad adicional de conocimiento cero, que garantiza que la prueba $\pi$ no revela información adicional sobre el \textit{witness} $w$ más allá de su existencia. Conceptualmente, podemos separar completamente la corrección del cómputo de la confidencialidad de los datos involucrados.

Para concluir con esta sección, hablaremos sobre el parámetro de seguridad $\securityparam$, que es una abstracción en sistemas criptográficos. Intuitivamente, $\securityparam$ mide el nivel de protección del sistema criptográfico frente a ataques. A mayor $\securityparam$, menor probabilidad de que un algoritmo adversario logre romper las garantías de seguridad del sistema. En la práctica, $\securityparam$ se relaciona en forma directa con el número de bits de seguridad del sistema: si decimos que ofrece $128$ bits de seguridad, significa que el mejor ataque conocido requiere del orden de $2^{128}$ operaciones para tener éxito, lo cual es computacionalmente inviable con la tecnología actual (y por eso $128$ suele ser el nivel deseado). Incrementar el valor de $\securityparam$ aumenta la dificultad de los ataques, pero también incrementa el tiempo de ejecución de los algoritmos, así como el tamaño de los argumentos y los parámetros públicos.

Como nota al pie, la notación $1^\securityparam$ utilizada como entrada del algoritmo $\Setup$ parece insensata, pero es una convención que indica que el parámetro se proporciona en forma unaria. En particular y como es indicado en \cite{Goldreich2001}, esto evita que algoritmos que dependen exponencialmente de $\securityparam$ aparezcan artificialmente como eficientes debido a una codificación binaria del parámetro, y que los algoritmos considerados eficientes sean precisamente aquellos cuyo tiempo de ejecución es polinomial en la longitud de sus entradas, preservando así la interpretación estándar.

En la siguiente sección veremos cómo representar cómputos generales mediante circuitos aritméticos, lo cual permitirá instanciar relaciones $\Relation$ adecuadas para la construcción de \SNARK concretos. Se remite al lector a \cite{DBLP:journals/ftsec/Thaler22}, donde podrá observar que hay múltiples formas de construir una \SNARK, muy diferentes a la que desarrollaremos en este capítulo.
